\color{uniblue}\section[Выбор уставок функции пуска охлаждения трансформатора]{\sloppy Выбор уставок функции пуска охлаждения \mbox{трансформатора}} 
\color{black}

\begin{enumerate}[label=\arabic{section}.\arabic*, labelsep=4pt, leftmargin=0pt, itemindent=57pt]

\item
Функция РТПО применяется на трансформаторах с высшим напряжением 35 кВ, для которых необходим внешний пуск охлаждения по току.
В общем случае ток срабатывания функции пуска охлаждения стороны высшего (низшего) напряжения трансформатора определяется по выражению:
\begin{equation}
    \label{eq:rtpo}
    I_{\textsl{ср}} = K_{\textsl{отс}}\cdot I_{\textsl{ном.ст.}},
    \end{equation}
    \begin{eqexpl}[15mm]
    \item{$I_{\textsl{ном.ст.}}$} номинальное значение тока стороны ВН (НН) трансформатора, А;
    \item{$K_{\textsl{отс}}$} коэффициент отстройки, принимается равным 0,95.
    \end{eqexpl}
\item
Выбор тока срабатывания для пуска охлаждения трансформатора по отстройке от номинального значения тока стороны трансформатора основан на \cite{rd3446501}, в случае наличия дополнительных требований к выбору уставки по току для пуска охлаждения от завода-изготовителя трансформатора, их также необходимо учитывать.

\end{enumerate}                              